\section*{Chapitre \ref{ch:b1:stratigraphic-scale} --- L'échelle stratigraphique}

\begin{solution}{exo:b1:stratigraphic-scale:1}
(a)~Chronostratigraphique, un \omterm{def:b1:stratigraphic-scale:chronostratigraphy}{système}. (b)~Lithostratigraphique, une \omterm{def:b1:stratigraphic-scale:lithostratigraphy}{formation}.
(c)~Géochronologique, une époque (sa roche est la \omterm{def:b1:stratigraphic-scale:chronostratigraphy}{Série} du Crétacé inférieur).
(d)~Biostratigraphique, une \omterm{def:b1:stratigraphic-scale:biozone}{biozone}. (e)~Chronostratigraphique, un \omterm{def:b1:stratigraphic-scale:chronostratigraphy}{étage}.
(f)~Géochronologique, une époque. (g)~Lithostratigraphique, un \omterm{def:b1:stratigraphic-scale:lithostratigraphy}{membre}.
\end{solution}

\begin{solution}{exo:b1:stratigraphic-scale:2}
Paléozoïque $538.8 - 251.902 = \qty{286.9}{\Ma}$, \qty{53.2}{\%} du
Phanérozoïque~; Mésozoïque $251.902 - 66.00 = \qty{185.9}{\Ma}$, \qty{34.5}{\%}~;
Cénozoïque \qty{66.0}{\Ma}, \qty{12.2}{\%}. Le Phanérozoïque fait
$538.8/4567 = \qty{11.8}{\%}$ du temps écoulé depuis la base de l'Hadéen.
\end{solution}

\begin{solution}{exo:b1:stratigraphic-scale:3}
A, B, C, D se sont déposés dans cet ordre (superposition), horizontalement (horizontalité
originelle), donc leur inclinaison est un événement ultérieur. Le dyke E coupe A à D, donc il est
plus jeune que D (recoupement)~; on ne peut dire d'après la description s'il est venu avant ou après
l'inclinaison. La surface d'érosion coupe D et E, donc elle est
plus jeune que les deux (recoupement). F et G se sont ensuite déposés dessus
(superposition). La faille décale A à F mais pas G~: elle a joué après F et avant G
(recoupement). Ordre~: A, B, C, D~; inclinaison et intrusion de E (ordre
indécidé)~; érosion~; F~; faille~; G.
\end{solution}

\begin{solution}{exo:b1:stratigraphic-scale:4}
$r = 30/(85.40 - 84.60) = \qty{37.5}{m}$ par million d'années. La \omterm{def:b1:stratigraphic-scale:biozone}{première apparition}
se situe $12/37.5 = \qty{0.32}{\Ma}$ après la cendre inférieure~: $85.40 - 0.32 =
\qty{85.08}{\Ma}$.
\end{solution}

\begin{solution}{exo:b1:stratigraphic-scale:5}
(a)~Quatre barres verticales en fonction de la hauteur. (b)~B et C se rencontrent ensemble de
\qtyrange{26}{50}{m}~: une \omterm{def:b1:stratigraphic-scale:biozone}{biozone} de co-extension épaisse de \qty{24}{m}. La \omterm{def:b1:stratigraphic-scale:biozone}{biozone
d'intervalle} de la \omterm{def:b1:stratigraphic-scale:biozone}{première apparition} de B à celle de D va de \qtyrange{12}{41}{m},
soit \qty{29}{m}. (c)~Dans la seconde coupe les \omterm{def:b1:stratigraphic-scale:biozone}{premières apparitions} de B (\qty{3}{m})
et de C (\qty{9}{m}) peuvent être placées, donc la base de la \omterm{def:b1:stratigraphic-scale:biozone}{biozone d'intervalle} de (b) et
la base de la \omterm{def:b1:stratigraphic-scale:biozone}{biozone} de co-extension sont reconnues~; son sommet (la \omterm{def:b1:stratigraphic-scale:biozone}{première
apparition} de D) et le sommet du recouvrement vu dans la coupe~1 ne peuvent être vérifiés,
D étant absent. Les \omterm{def:b1:stratigraphic-scale:biozone}{premières apparitions} de B et de C ($\qty{12}{m} \leftrightarrow
\qty{3}{m}$, $\qty{26}{m} \leftrightarrow \qty{9}{m}$) sont les lignes les plus sûres~; les
\omterm{def:b1:stratigraphic-scale:biozone}{dernières apparitions} (B à \qty{50}{m} et \qty{15}{m}) dépendent davantage de la préservation.
\end{solution}

\begin{solution}{exo:b1:stratigraphic-scale:6}
(a)~$60/20 = \qty{3}{\Ma}$~: le sommet du grès est plus jeune de trois millions d'années
dans le puits du côté des terres. (b)~Un million d'années de boue à \qty{30}{m} par
million d'années~: la cendre se trouve à \qty{30}{m} au-dessus du sommet du grès.
(c)~Quand la cendre est tombée, le rivage avait avancé de \qty{20}{km} vers l'intérieur au-delà du puits
du côté du large, donc il se trouvait encore à \qty{40}{km} au large du puits du côté des terres, qui était
en terre ferme. La cendre est tombée sur la surface terrestre en ce lieu, sous le grès que la
mer a déposé deux millions d'années plus tard --- si elle n'a pas d'abord été emportée,
comme une couche meuble sur terre peut fort bien l'être. Le même banc de cendre est au-dessus de la limite de \omterm{def:b1:stratigraphic-scale:lithostratigraphy}{formation} dans
un puits et en dessous dans l'autre~: la limite traverse le temps.
\end{solution}

\begin{solution}{exo:b1:stratigraphic-scale:7}
(a)~Lochkovien $1.36/419.62 = \qty{0.32}{\%}$~; Indusien $0.024/251.902 =
\qty{0.0095}{\%}$, environ 34 fois plus petit. La base de l'Indusien se trouve parmi des bancs de cendre
datés avec précision dans la coupe du GSSP elle-même~; une grande incertitude suggère une
limite dont l'âge doit être transporté depuis des niveaux datés plus lointains.
(b)~$251.941 - 251.880 = \qty{0.061}{\Ma} = 61$ milliers d'années, avec
l'incertitude $\sqrt{37^2 + 31^2} = 48$ milliers d'années~: les
\num[uncertainty-mode=separate]{60 \pm 48} milliers d'années publiés.
\end{solution}

\begin{solution}{exo:b1:stratigraphic-scale:8}
(a)~$m = 3000 \times \frac43\pi\,(5000)^3 = \qty{1.57e15}{kg}$~; la poussière,
$60m = \qty{9.42e16}{kg}$. (b)~La surface de la Terre est $4\pi(6.378\times10^6)^2
= \qty{5.11e14}{m^2}$~: \qty{184}{kg} par mètre carré. (c)~$184/2500 =
\qty{0.074}{m}$, environ \qty{7}{cm}. (d)~L'épaisseur varie comme le cube du
diamètre~: \qty{1.6}{cm} pour \qty{6}{km}, \qty{20}{cm} pour \qty{14}{km}. L'ordre
de grandeur, des centimètres aux décimètres, est celui de l'argile de limite,
épaisse de \qty{50}{cm} à El Kef~; la comparaison ne peut être que grossière, puisque la poussière
n'a pu tomber uniformément et qu'une argile recueille aussi le sédiment ordinaire de
son temps.
\end{solution}

\begin{solution}{exo:b1:stratigraphic-scale:9}
(a)~$\num{60e3} \times \num{2e-6} = \qty{0.12}{m}$, \qty{12}{cm}~; pour 12 et 108
milliers d'années, \qty{2.4}{cm} et \qty{21.6}{cm}. (b)~$\num{32e3} \times
\num{10e-6} = \qty{0.32}{m}$. (c)~Une extinction occupe des centimètres à des décimètres
de roche~: il faut prélever les échantillons centimètre par centimètre, et un seul échantillon
de \qty{10}{cm} d'épaisseur dans la coupe condensée fait la moyenne de cinquante mille ans ---
l'essentiel de l'événement.
\end{solution}

\begin{solution}{exo:b1:stratigraphic-scale:10}
(a)~$u = R/s = 10 \times 2000 = \qty{20000}{m}$ par million d'années,
soit \qty{20}{km} par million d'années~; sur \qty{30}{km}, $\Delta t = \qty{1.5}{\Ma}$.
(b)~La plaine est maintenant la surface $z = s x + a t$. Le rivage vérifie
$s x + a t = z_0 + R t$, donc $x(t) = (z_0 + (R - a)t)/s$ et
$\mathrm dx/\mathrm dt = (R - a)/s$, positif (vers l'intérieur) tant que $a < R$. Quand
$a > R$ la vitesse est négative~: le rivage se déplace vers le large, c'est une régression, et
le grès de rivage est alors plus jeune vers le large. (c)~Une falaise longue de \qty{200}{m} couvre,
à \qty{20}{km} par million d'années, $0.2/20 = \qty{0.01}{\Ma}$ --- dix mille
ans de diachronisme, bien en dessous de ce que les fossiles ou la plupart des datations peuvent résoudre~; la roche de
part et d'autre du contact est la même partout le long de lui.
\end{solution}

\begin{solution}{exo:b1:stratigraphic-scale:11}
(a)~$f = (25 - 10)/(50 - 10) = 0.375$ et $t = 100.50 - 0.375 \times 0.80 =
\qty{100.20}{\Ma}$, avec
\[
  \sigma_t = \sqrt{(0.625 \times 0.05)^2 + (0.375 \times 0.03)^2} = \qty{0.033}{\Ma}.
\] (b)~$g(f) = (1 - f)^2\sigma_1^2 + f^2\sigma_2^2$
a pour dérivée $g'(f) = -2(1 - f)\sigma_1^2 + 2f\sigma_2^2$, nulle en $f^* =
\sigma_1^2/(\sigma_1^2 + \sigma_2^2) = 0.0025/0.0034 = 0.735$, plus près du banc le mieux
daté~; là $\sigma_t = \sigma_1\sigma_2/\sqrt{\sigma_1^2 + \sigma_2^2} =
\qty{0.026}{\Ma}$, plus petit que chacun des deux. (c)~$g$ est une parabole ouverte vers le haut, donc
sur $[0, 1]$ elle est la plus grande à une extrémité~: $g(0) = \sigma_1^2$, $g(1) = \sigma_2^2$.
D'où $\sigma_t \le \max(\sigma_1, \sigma_2)$ entre les bancs.
\end{solution}

\begin{solution}{exo:b1:stratigraphic-scale:12}
(a)~La limite est à mi-chemin en profondeur~: $f = 0.5$, \qty{29.50}{\Ma}. (b)~Avec
$\varphi_0\lambda = 0.63 \times 1960 = \qty{1234.8}{m}$, l'épaisseur solide de
\qtyrange{600}{1000}{m} est $400 - 1234.8\,(e^{-600/1960} - e^{-1000/1960}) =
\qty{232.2}{m}$, et de \qtyrange{200}{600}{m}, \qty{194.2}{m}~; au total
\qty{426.3}{m}. La fraction des grains déposés entre la cendre inférieure et la
limite est $232.2/426.3 = 0.545$, et l'âge est $30.0 - 0.545 =
\qty{29.46}{\Ma}$. (c)~La seconde estimation est plus jeune de \num{45} milliers d'années.
C'est la meilleure~: la moitié profonde a perdu plus de son eau, donc elle contient
plus de grains, et plus de temps, par mètre que la moitié moins profonde (porosité moyenne
\num{0.42} contre \num{0.51}).
\end{solution}

\subsection*{Problème~: Corréler trois coupes à travers un bassin}

\begin{solution}{pb:b1:stratigraphic-scale:1}
\begin{enumerate}
  \item Trois colonnes à, disons, \qty{1}{cm} pour \qty{10}{m}, reliées par les bancs de cendre
        et par les limites de zones~; le coin du grès de Kell s'amincit de
        \qty{45}{m} en A à \qty{6}{m} en C.
  \item Zone~2~: A $65.8 - 36.6 = \qty{29.2}{m}$, B \qty{20.8}{m}, C
        \qty{10.0}{m}. Zone~3 levée~: A \qty{14.2}{m}, B \qty{14.0}{m}, C
        \qty{7.2}{m}.
  \item La zone~1 commence à la \omterm{def:b1:stratigraphic-scale:biozone}{première apparition} de \emph{V.~primum}, qui se trouve
        depuis la base de chaque coupe~: sa \omterm{def:b1:stratigraphic-scale:biozone}{première apparition} se situe en dessous d'elles.
  \item Il ne vit que dans le grès~: il suit le faciès, et le faciès
        est d'âges différents d'un lieu à l'autre. En A il se trouve dans les zones~1
        et~2, en B et C seulement dans la zone~1.
  \item Dans la zone~2 dans les trois (A \qty{49.0}{m}, B \qty{40.0}{m}, C
        \qty{19.0}{m}). Elle est cohérente d'une coupe à l'autre ici, et pourrait donc
        subdiviser la zone~2, mais une \omterm{def:b1:stratigraphic-scale:biozone}{dernière apparition} dépend de la préservation~;
        il faudrait d'abord la vérifier dans davantage de coupes.
  \item Les lignes joignent \qty{36.6}{m}, \qty{31.2}{m} et \qty{14.8}{m}, puis
        \qty{65.8}{m}, \qty{52.0}{m} et \qty{24.8}{m}. La zone~2 s'amincit dans le
        rapport $29.2 : 20.8 : 10.0$, environ $2.9 : 2.1 : 1$~: si la zone a duré
        le même temps partout, le sédiment s'est accumulé environ trois fois plus
        vite en A qu'en C.
  \item A~: zone~2 ($45.0 > 36.6$)~; B~: zone~1 ($20.0 < 31.2$)~; C~: zone~1
        ($6.0 < 14.8$).
  \item Non~: il est en zone~2 en A, en zone~1 en B et C, donc il est plus jeune à
        l'ouest. Les argilites marines se sont étendues vers l'ouest sur le sable de rivage~: le
        rivage s'est déplacé vers l'ouest, vers l'intérieur, dans une transgression.
  \item Une cendre tombe sur tout le bassin en un instant géologique et se dépose
        sur n'importe quel sédiment~; le sommet du grès suit le rivage mobile
        et est \omterm{def:b1:stratigraphic-scale:diachronous}{diachrone}.
  \item La continuité latérale (le banc se poursuit à travers le bassin), et
        la superposition (cendre~1 sous cendre~2 dans chaque coupe, ce qui vérifie la
        \omterm{def:b1:cross-sections:log}{corrélation}).
  \item A $(71.0 - 5.0)/1.20 = \qty{55.0}{m}$ par million d'années~; B
        $48.0/1.20 = \qty{40.0}{}$~; C $24.0/1.20 = \qty{20.0}{}$.
  \item A est la plus proche de la terre, où le sédiment était fourni~; moins de celui-ci
        a atteint les eaux plus profondes à l'est.
  \item $0.1/55 = \num{1.8}$ milliers d'années en A, $\num{2.5}$ en B, $\num{5.0}$
        en C~: A résout le temps le plus finement.
  \item A $152.40 - 40.0/55.0 = \qty{151.67}{\Ma}$~; B $152.40 - 12.0/40.0 =
        \qty{152.10}{\Ma}$~; C $152.40 - 3.0/20.0 = \qty{152.25}{\Ma}$.
  \item C à B~: $20/(152.25 - 152.10) = \qty{133}{km}$ par million d'années~; B à A~:
        $25/(152.10 - 151.67) = \qty{58.5}{km}$ par million d'années. Elle s'est
        ralentie.
  \item A $152.40 - 60.8/55.0 = \qty{151.29}{\Ma}$~; B $152.40 - 44.0/40.0 =
        \qty{151.30}{\Ma}$~; C $152.40 - 21.8/20.0 = \qty{151.31}{\Ma}$.
  \item Ils couvrent \num{16} milliers d'années, moins que l'incertitude de chacun
        (environ \num{28} milliers d'années, question~19)~: la \omterm{def:b1:stratigraphic-scale:biozone}{première apparition} est
        synchrone dans le bassin dans la résolution des datations, ce qui
        en fait un bon repère.
  \item A $152.40 - 31.6/55.0 = \qty{151.826}{\Ma}$~; B $152.40 - 23.2/40.0 =
        \qty{151.820}{\Ma}$~; C $152.40 - 11.8/20.0 = \qty{151.810}{\Ma}$.
  \item $f = 0.479$, $0.483$, $0.492$~; $\sigma_t = \sqrt{(1 - f)^2 + f^2} \times
        0.04 = \qty{0.028}{\Ma}$ dans chacune, \num{28} milliers d'années.
  \item Moyenne \qty{151.819}{\Ma}~; dispersion $151.826 - 151.810 = \num{16}$ milliers
        d'années, plus petite que l'incertitude de chaque âge~: les trois s'accordent.
  \item L'interpolation suppose une seule vitesse de la cendre~1 à la cendre~2. En A cet
        intervalle contient du grès de rivage jusqu'à \qty{45}{m} et de l'argilite au-dessus~;
        le sable et la boue n'ont pas à s'être accumulés à la même vitesse, et ils
        se compactent différemment (\cref{ch:b1:diagenesis-basins}). Si le sable
        s'est accumulé plus vite, la partie gréseuse contient moins de temps que son
        épaisseur ne le suggère, et la \omterm{def:b1:stratigraphic-scale:biozone}{première apparition} est plus ancienne que calculée
        --- l'âge de A est bien le plus ancien des trois.
  \item B~: la \omterm{def:b1:stratigraphic-scale:biozone}{première apparition} se situe dans l'argilite marine, \qty{11}{m} au-dessus du
        sommet du grès, sans changement de faciès près du niveau, et la
        vitesse donne \num{2.5} milliers d'années par \qty{10}{cm}. A résout le temps
        mieux mais place le niveau dans le grès de rivage, \qty{8.4}{m}
        sous un changement de faciès~; C est condensée, \num{5} milliers d'années par
        \qty{10}{cm}. C est la pire pour la résolution, A pour le faciès.
  \item À partir des âges moyens des deux \omterm{def:b1:stratigraphic-scale:biozone}{premières apparitions}, $151.819 - 151.302 =
        \qty{0.52}{\Ma}$~: environ un demi-million d'années.
  \item La limite est de \qty[uncertainty-mode=separate]{151.82(3)}{\Ma}.
\end{enumerate}
\end{solution}
